\section{Introduction}
\label{sec:intro}

Phishing remains one of the most prevalent and costly cybersecurity threats, and
automated detectors deployed across web browsers, mail gateways, and enterprise
security pipelines are a primary line of defense. A phishing verdict is also a consequential
act: marking a page as phishing drives a blocklist entry, a browser warning, or a
take-down request that pulls a live service off the web. Large language models have sharply raised the accuracy of
such detectors, reading a page's URL, Document Object Model (DOM), rendered screenshot, and certificate
to decide whether it impersonates a brand it does not own~\cite{tan2024phishllm,
li2024knowphish}. Yet accuracy is only half of what a deployment needs: the
operator who must justify a take-down and the user who sees a warning both need a
\emph{reason} they can check, and the system needs a criterion for \emph{when} a
verdict is trustworthy enough to act on without review. For a decision that can be
contested, the bottleneck is no longer whether the model is usually right, but
whether a particular verdict can be \emph{trusted and explained}.

Existing detectors approach this only indirectly. One line runs a single LLM over
the page and reports a class probability or a verbalized
confidence~\cite{tan2024phishllm}; reliability is then a property of that one
model's calibration~\cite{guo2017calibration}. A second line aggregates several
models or several samples -- self-consistency over one
model~\cite{wang2023selfconsistency}, majority voting over an
ensemble~\cite{ensembleornot2024}, or weighted label agreement in the style of
crowd-sourcing and label-free evaluation~\cite{dawid1979maximum,tci2025} -- and
reads reliability off how often the votes concur. A third, recent line builds
multi-agent detectors whose agents specialize on different signals and whose
rationales are merged into a single explanation~\cite{phishdebate2025,
multiphishguard2025}. Together these families cover much of the space of accuracy
and aggregation.

\begin{figure*}[!h]
\centering
% ---- palette -------------------------------------------------------------
\definecolor{PPInk}{HTML}{243443}
\definecolor{PPBlue}{HTML}{DCEFF7}
\definecolor{PPMint}{HTML}{DCF2D8}
\definecolor{PPCream}{HTML}{FFF2C6}
\definecolor{PPPeach}{HTML}{FCE2D4}
\definecolor{PPLav}{HTML}{E9E6FF}
\definecolor{PPGreen}{HTML}{2F8F5B}
\definecolor{PPRed}{HTML}{B24A4A}
\definecolor{PPGold}{HTML}{D49B22}
% agent-identity hues (kept distinct from the green/amber/red verdict colors)
\definecolor{PPA}{HTML}{2C6FB0}   % A1 text    -- blue
\definecolor{PPB}{HTML}{6E54C8}   % A2 vision  -- violet
\definecolor{PPC}{HTML}{1F8E8B}   % A3 security-- teal
% small inline half-disc used as the "partial" mark (matches the board stamp)
\newcommand{\hpart}{\tikz[baseline=-0.42ex]{\begin{scope}\clip (0,0) circle (0.92mm);
  \fill[PPGold!85!black] (-1mm,-1mm) rectangle (0,1mm);\end{scope}
  \draw[line width=0.25pt, PPGold!85!black] (0,0) circle (0.92mm);}}
\resizebox{\textwidth}{!}{%
\begin{tikzpicture}[
  font=\sffamily\scriptsize,
  >=Latex,
  stage/.style={font=\sffamily\tiny\bfseries, text=PPInk, align=left},
  sbadge/.style={circle, fill=PPInk, text=white, font=\sffamily\tiny\bfseries,
    minimum size=3.4mm, inner sep=0pt},
  card/.style={draw=PPInk!45, line width=0.35pt, rounded corners=2.5pt,
    fill=white, inner sep=3pt, align=center,
    drop shadow={shadow xshift=0.25pt, shadow yshift=-0.25pt, fill=black!9}},
  agent/.style={card, fill=white, text width=26mm, minimum height=11.5mm, align=left},
  board/.style={card, fill=PPMint!22, minimum width=44mm, minimum height=41mm},
  gauge/.style={card, fill=PPLav!26, minimum width=33mm, minimum height=37mm},
  decision/.style={card, text width=34mm, minimum height=19mm, align=left},
  erow/.style={draw=PPInk!22, line width=0.3pt, rounded corners=2pt,
    fill=white, minimum width=40mm, minimum height=9.4mm, inner sep=0pt},
  toolchip/.style={draw=PPInk!28, line width=0.25pt, rounded corners=1pt,
    fill=PPInk!4, inner sep=1.3pt, font=\sffamily\fontsize{4.5}{5}\selectfont, text=PPInk!82},
  flow/.style={-{Latex[length=1.6mm]}, line width=0.55pt, draw=PPInk!72},
  weakflow/.style={-{Latex[length=1.35mm]}, line width=0.5pt, draw=PPRed!75, dashed},
  thin/.style={line width=0.25pt, draw=PPInk!30},
  pin/.style={circle, draw=white, line width=0.2pt, minimum size=2.7mm,
    inner sep=0pt, font=\sffamily\fontsize{4.2}{4.4}\selectfont\bfseries, text=white},
  stamp/.style={circle, line width=0.3pt, minimum size=4mm, inner sep=0pt,
    font=\sffamily\tiny\bfseries},
  % ---------- icon library (drawn, so they scale with the figure) ----------
  pics/ic text/.style={code={
    \draw[line width=0.3pt, rounded corners=0.2mm] (-1mm,-1.3mm) rectangle (1mm,1.3mm);
    \draw[line width=0.28pt] (-0.6mm,0.8mm)--(0.6mm,0.8mm);
    \draw[line width=0.28pt] (-0.6mm,0.3mm)--(0.6mm,0.3mm);
    \draw[line width=0.28pt] (-0.6mm,-0.2mm)--(0.6mm,-0.2mm);
    \draw[line width=0.28pt] (-0.6mm,-0.7mm)--(0.2mm,-0.7mm);}},
  pics/ic eye/.style={code={
    \draw[line width=0.32pt] (-1.25mm,0) .. controls (-0.5mm,1.05mm) and (0.5mm,1.05mm) .. (1.25mm,0)
       .. controls (0.5mm,-1.05mm) and (-0.5mm,-1.05mm) .. (-1.25mm,0) -- cycle;
    \fill (0,0) circle (0.46mm);}},
  pics/ic shield/.style={code={
    \draw[line width=0.32pt, line join=round] (0,1.45mm) -- (1.2mm,0.85mm) -- (1.2mm,-0.35mm)
      .. controls (1.2mm,-1.05mm) and (0.55mm,-1.4mm) .. (0,-1.6mm)
      .. controls (-0.55mm,-1.4mm) and (-1.2mm,-1.05mm) .. (-1.2mm,-0.35mm) -- (-1.2mm,0.85mm) -- cycle;
    \draw[line width=0.34pt, line cap=round] (-0.5mm,-0.05mm)--(-0.08mm,-0.5mm)--(0.6mm,0.5mm);}},
  pics/ic lock/.style={code={
    \draw[line width=0.3pt, rounded corners=0.3mm] (-0.95mm,-1.15mm) rectangle (0.95mm,0.25mm);
    \draw[line width=0.3pt] (-0.5mm,0.25mm) .. controls (-0.5mm,1.05mm) and (0.5mm,1.05mm) .. (0.5mm,0.25mm);
    \fill (0,-0.45mm) circle (0.24mm);}},
  pics/ic bank/.style={code={
    \draw[line width=0.3pt, line join=round] (-1.35mm,0.3mm)--(0,1.25mm)--(1.35mm,0.3mm);
    \draw[line width=0.3pt] (-1.15mm,0.3mm)--(1.15mm,0.3mm);
    \foreach \cx in {-0.92,-0.31,0.31,0.92} \draw[line width=0.34pt] (\cx mm,0.18mm)--(\cx mm,-0.95mm);
    \draw[line width=0.34pt] (-1.25mm,-0.98mm)--(1.25mm,-0.98mm);}},
  pics/ic warn/.style={code={
    \draw[line width=0.3pt, line join=round] (0,1.3mm)--(1.3mm,-1mm)--(-1.3mm,-1mm)--cycle;
    \draw[line width=0.36pt, line cap=round] (0,0.65mm)--(0,-0.2mm);
    \fill (0,-0.58mm) circle (0.2mm);}},
  pics/ic packet/.style={code={
    \draw[line width=0.3pt, fill=white] (-1.05mm,-1.45mm) -- (0.55mm,-1.45mm) -- (1.05mm,-0.95mm)
       -- (1.05mm,1.45mm) -- (-1.05mm,1.45mm) -- cycle;
    \draw[line width=0.26pt] (0.55mm,-1.45mm) -- (0.55mm,-0.95mm) -- (1.05mm,-0.95mm);
    \draw[line width=0.26pt] (-0.65mm,0.85mm)--(0.65mm,0.85mm);
    \draw[line width=0.26pt] (-0.65mm,0.35mm)--(0.65mm,0.35mm);
    \draw[line width=0.26pt] (-0.65mm,-0.15mm)--(0.65mm,-0.15mm);}},
  pics/ic half/.style={code={
    \begin{scope} \clip (0,0) circle (1.25mm); \fill (-1.3mm,-1.3mm) rectangle (0,1.3mm); \end{scope}
    \draw[line width=0.3pt] (0,0) circle (1.25mm);}},
  pics/ic pause/.style={code={
    \draw[line width=0.55pt, line cap=round] (-0.5mm,1mm)--(-0.5mm,-1mm);
    \draw[line width=0.55pt, line cap=round] (0.5mm,1mm)--(0.5mm,-1mm);}}
]

% ===== stage header band =================================================
\foreach \x/\n/\t in {0/1/{page under test}, 3.6/2/{agent perspectives},
   7.4/3/{shared evidence board}, 11.4/4/{grounded agreement}, 15.6/5/{outcome}} {
  \node[sbadge] (sb\n) at (\x-0.98,2.66) {\n};
  \node[stage, anchor=west] at (\x-0.74,2.66) {\t};
}
\foreach \a/\b in {1.95/2.32, 5.78/6.12, 9.62/9.98, 13.82/14.18} {
  \draw[thin, -{Latex[length=1.1mm]}, draw=PPInk!22] (\a,2.66)--(\b,2.66); }

% ===== (1) page under test : detailed browser ============================
\begin{scope}[shift={(0,0.2)}]
  \draw[card, fill=PPBlue!16] (-1.5,-1.55) rectangle (1.5,1.95);
  \fill[PPInk!9, rounded corners=2.5pt] (-1.5,1.55) rectangle (1.5,1.95);
  \draw[thin] (-1.5,1.55) -- (1.5,1.55);
  \fill[PPRed!60]   (-1.34,1.75) circle (0.052);
  \fill[PPGold!75]  (-1.20,1.75) circle (0.052);
  \fill[PPGreen!60] (-1.06,1.75) circle (0.052);
  % address bar with warning + deceptive domain
  \draw[thin, fill=white, rounded corners=1.4pt] (-1.32,1.12) rectangle (1.32,1.42);
  \pic[PPRed] at (-1.16,1.27) {ic warn};
  \node[anchor=west, font=\sffamily\fontsize{5}{5.4}\selectfont, text=PPInk!80]
    at (-1.0,1.27) {secure-bank-login.example};
  % brand lock-up (bank glyph, not a trust shield)
  \pic[PPInk!75] at (-0.78,0.74) {ic bank};
  \node[anchor=west, font=\sffamily\fontsize{6.4}{6.8}\selectfont\bfseries, text=PPInk]
    at (-0.58,0.76) {BankSecure};
  \node[font=\sffamily\fontsize{4.4}{4.8}\selectfont, text=PPInk!55] at (0,0.42) {Online Banking Login};
  % credential fields
  \draw[thin, fill=PPInk!3, rounded corners=1pt] (-1.0,0.05) rectangle (1.0,0.27);
  \node[anchor=west, font=\sffamily\fontsize{4.4}{4.8}\selectfont, text=PPInk!50] at (-0.92,0.16) {Username};
  \draw[thin, fill=PPInk!3, rounded corners=1pt] (-1.0,-0.28) rectangle (1.0,-0.06);
  \node[anchor=west, font=\sffamily\fontsize{4.4}{4.8}\selectfont, text=PPInk!50] at (-0.92,-0.17) {Password};
  % sign-in button
  \draw[draw=none, fill=PPInk!72, rounded corners=1.2pt] (-0.62,-0.66) rectangle (0.62,-0.42);
  \node[font=\sffamily\fontsize{4.8}{5.2}\selectfont\bfseries, text=white] at (0,-0.54) {Sign in};
  % fake trust seal
  \pic[PPGreen!70] at (-0.78,-1.05) {ic lock};
  \node[anchor=west, font=\sffamily\fontsize{4}{4.4}\selectfont, text=PPInk!45] at (-0.62,-1.05)
    {256-bit secure};
\end{scope}
\node[font=\sffamily\fontsize{5}{5.4}\selectfont\bfseries, text=PPInk] at (0,2.0) {suspicious login page};

% raw-signal card under the browser
\node[card, fill=white, text width=27mm, minimum height=8mm, align=center,
  font=\sffamily\fontsize{4.8}{6}\selectfont, text=PPInk!78] (sig) at (0,-2.05)
  {\textbf{raw page signals}\\ URL\,\textbullet\,DOM form\\ screenshot\,\textbullet\,logo};

% ===== (2) agent perspectives ============================================
\node[agent] (a1) at (3.6,1.18) {};
\node[agent] (a2) at (3.6,0.10) {};
\node[agent] (a3) at (3.6,-0.98) {};
\foreach \a/\col/\ic/\nm/\rd in {%
   a1/PPA/{ic text}/{A1\, text agent}/{URL \textbullet\ DOM},
   a2/PPB/{ic eye}/{A2\, vision agent}/{screenshot \textbullet\ logo},
   a3/PPC/{ic text}/{A3\, text agent}/{DOM \textbullet\ form}} {
  \fill[\col, rounded corners=1pt] ($(\a.north west)+(0.6mm,-0.6mm)$) rectangle ($(\a.south west)+(1.7mm,0.6mm)$);
  \node[circle, draw=\col!65, line width=0.3pt, fill=\col!12, minimum size=4.6mm, inner sep=0pt]
    (\a-av) at ($(\a.west)+(4.2mm,0)$) {};
  \pic[\col] at (\a-av.center) {\ic};
  \node[anchor=west, font=\sffamily\fontsize{5.6}{6}\selectfont\bfseries, text=PPInk]
    at ($(\a-av.east)+(0.7mm,1.3mm)$) {\nm};
  \node[anchor=west, font=\sffamily\fontsize{4.6}{5}\selectfont, text=PPInk!58]
    at ($(\a-av.east)+(0.7mm,-1.3mm)$) {reads \rd};
}
\foreach \a in {a1,a2,a3} { \draw[flow] (1.55,0.2) -- (\a.west); }

% ===== (3) shared evidence board =========================================
\node[board] (board) at (7.4,0.05) {};
\node[font=\sffamily\fontsize{5.2}{5.6}\selectfont\bfseries, text=PPInk] at ($(board.north)+(0,-0.22)$)
  {cited cues with provenance \& page checks};

\node[erow] (e1) at ($(board.center)+(0,0.92)$) {};
\node[erow] (e2) at ($(board.center)+(0,-0.02)$) {};
\node[erow] (e3) at ($(board.center)+(0,-0.96)$) {};

% row 1 : logo -- partial
\node[pin, fill=PPA] at ($(e1.west)+(2.4mm,0)$) {1};
\node[pin, fill=PPB] at ($(e1.west)+(5.0mm,0)$) {2};
\node[anchor=west, font=\sffamily\fontsize{5.2}{5.6}\selectfont\bfseries, text=PPInk]
  at ($(e1.west)+(7.2mm,1.4mm)$) {logo resembles bank};
\node[anchor=west, font=\sffamily\fontsize{4.3}{4.7}\selectfont, text=PPInk!55]
  at ($(e1.west)+(7.2mm,-1.5mm)$) {visual brand cue};
\node[toolchip, anchor=east] at ($(e1.east)+(-5.4mm,0)$) {logo\,matcher};
\node[stamp, draw=PPGold, fill=PPGold!12, anchor=east] (st1) at ($(e1.east)+(-1.5mm,0)$) {};
\pic[PPGold!85!black, scale=0.7] at (st1.center) {ic half};

% row 2 : form off-brand -- confirmed
\node[pin, fill=PPA] at ($(e2.west)+(2.4mm,0)$) {1};
\node[pin, fill=PPC] at ($(e2.west)+(5.0mm,0)$) {3};
\node[anchor=west, font=\sffamily\fontsize{5.2}{5.6}\selectfont\bfseries, text=PPInk]
  at ($(e2.west)+(7.2mm,1.4mm)$) {form posts off-brand};
\node[anchor=west, font=\sffamily\fontsize{4.3}{4.7}\selectfont, text=PPInk!55]
  at ($(e2.west)+(7.2mm,-1.5mm)$) {form-action domain};
\node[toolchip, anchor=east] at ($(e2.east)+(-5.4mm,0)$) {DOM\,parser};
\node[stamp, draw=PPGreen, fill=PPGreen!12, text=PPGreen!75!black, anchor=east] at ($(e2.east)+(-1.5mm,0)$) {\cmark};

% row 3 : extra brand claim -- refuted by the brand detector
\node[pin, fill=PPC] at ($(e3.west)+(2.4mm,0)$) {3};
\node[anchor=west, font=\sffamily\fontsize{5.2}{5.6}\selectfont\bfseries, text=PPInk]
  at ($(e3.west)+(7.2mm,1.4mm)$) {extra brand claimed};
\node[anchor=west, font=\sffamily\fontsize{4.3}{4.7}\selectfont, text=PPRed!70]
  at ($(e3.west)+(7.2mm,-1.5mm)$) {not shown on page};
\node[toolchip, anchor=east] at ($(e3.east)+(-5.4mm,0)$) {brand\,det.};
\node[stamp, draw=PPRed, fill=PPRed!10, text=PPRed!85, anchor=east] at ($(e3.east)+(-1.5mm,0)$) {\xmark};

% legend strip (two compact centered lines that fit inside the board)
\node[anchor=center, align=center, text width=41mm,
  font=\sffamily\fontsize{4.3}{5.4}\selectfont, text=PPInk!64] at ($(board.south)+(0,0.34)$)
  {\textbf{pins}\,=\,which agents cited the cue\\[0.2pt]
   \textbf{stamp}\,=\,page re-check:\ \ \textcolor{PPGreen!70!black}{\cmark}\,confirm \ \ \hpart\,partial \ \ \textcolor{PPRed!85}{\xmark}\,refute};

\foreach \a in {a1,a2,a3} { \draw[flow] (\a.east) -- (board.west); }

% ===== (4) grounded agreement gauge ======================================
\node[gauge] (g) at (11.4,0.05) {};
\node[font=\sffamily\fontsize{5.2}{5.6}\selectfont\bfseries, text=PPInk] at ($(g.north)+(0,-0.22)$)
  {trust = agreement $\times$ verify};
% shared+verified meter (high)
\node[font=\sffamily\fontsize{4.7}{5}\selectfont\bfseries, text=PPInk, anchor=west] at ($(g.center)+(-1.32,0.86)$)
  {shared + verified};
\draw[thin, fill=white, rounded corners=1pt] ($(g.center)+(-1.32,0.42)$) rectangle ($(g.center)+(1.32,0.66)$);
\fill[PPGreen!72, rounded corners=1pt] ($(g.center)+(-1.32,0.42)$) rectangle ($(g.center)+(0.78,0.66)$);
\node[font=\sffamily\fontsize{4.4}{4.8}\selectfont, text=PPGreen!55!black, anchor=west] at ($(g.center)+(-1.32,0.20)$)
  {trust rises $\rightarrow$ act};
% label-only meter (low)
\node[font=\sffamily\fontsize{4.7}{5}\selectfont\bfseries, text=PPInk, anchor=west] at ($(g.center)+(-1.32,-0.42)$)
  {label-only agreement};
\draw[thin, fill=white, rounded corners=1pt] ($(g.center)+(-1.32,-0.86)$) rectangle ($(g.center)+(1.32,-0.62)$);
\fill[PPRed!60, rounded corners=1pt] ($(g.center)+(-1.32,-0.86)$) rectangle ($(g.center)+(-0.66,-0.62)$);
\node[font=\sffamily\fontsize{4.4}{4.8}\selectfont, text=PPRed!70!black, anchor=west] at ($(g.center)+(-1.32,-1.08)$)
  {stays low $\rightarrow$ abstain};
% threshold line tau across both meters
\draw[line width=0.4pt, dash pattern=on 0.7mm off 0.5mm, draw=PPInk!75]
  ($(g.center)+(0.30,0.80)$) -- ($(g.center)+(0.30,-0.98)$);
\node[font=\sffamily\fontsize{4.8}{5}\selectfont, text=PPInk!80] at ($(g.center)+(0.30,0.98)$) {$\tau$};
\node[font=\sffamily\fontsize{4}{4.4}\selectfont, text=PPInk!50, align=center] at ($(g.center)+(0,-1.42)$)
  {$\tau$ set on a labeled split};
\draw[flow] (board.east) -- (g.west);

% ===== (5) outcome =======================================================
\node[decision, fill=PPMint!42, minimum height=20mm] (act) at (15.6,1.05) {};
\node[decision, fill=PPPeach!48, minimum height=20mm] (abstain) at (15.6,-1.05) {};
% --- act ---
\node[circle, draw=PPGreen!70, fill=PPGreen!12, minimum size=5mm, inner sep=0pt]
  (actic) at ($(act.north west)+(4mm,-4mm)$) {};
\pic[PPGreen!70!black] at (actic.center) {ic shield};
\node[anchor=west, font=\sffamily\fontsize{6.4}{6.8}\selectfont\bfseries, text=PPInk]
  at ($(actic.east)+(0.9mm,0)$) {act};
\node[anchor=north west, font=\sffamily\fontsize{4.8}{5.6}\selectfont, text=PPInk!80, text width=30mm, align=left]
  at ($(act.north west)+(2mm,-8.6mm)$) {block or warn, and return the verified \textbf{evidence packet}};
\pic[PPInk!72] at ($(act.south west)+(4.3mm,3mm)$) {ic packet};
\node[anchor=west, font=\sffamily\fontsize{4.4}{4.8}\selectfont, text=PPInk!72]
  at ($(act.south west)+(6.6mm,3mm)$) {form off-brand\,\textcolor{PPGreen!70!black}{\cmark}\ \ logo\,\hpart};
% --- abstain ---
\node[circle, draw=PPGold!75, fill=PPGold!14, minimum size=5mm, inner sep=0pt]
  (absic) at ($(abstain.north west)+(4mm,-4mm)$) {};
\pic[PPGold!80!black] at (absic.center) {ic pause};
\node[anchor=west, font=\sffamily\fontsize{6.4}{6.8}\selectfont\bfseries, text=PPInk]
  at ($(absic.east)+(0.9mm,0)$) {abstain};
\node[anchor=north west, font=\sffamily\fontsize{4.8}{5.6}\selectfont, text=PPInk!80, text width=30mm, align=left]
  at ($(abstain.north west)+(2mm,-8.6mm)$) {label may be right, but the reasons split or fail page checks};
\node[anchor=south west, font=\sffamily\fontsize{4.8}{5.4}\selectfont\bfseries, text=PPInk!82, align=left]
  at ($(abstain.south west)+(2mm,2.2mm)$) {$\rightarrow$ escalate to a human\\or deeper analysis};
\draw[flow] ($(g.east)+(0,0.66)$) -- (act.west);
\draw[weakflow] ($(g.east)+(0,-0.56)$) -- (abstain.west);

% ===== footer ============================================================
\node[font=\sffamily\fontsize{5.2}{5.6}\selectfont, text=PPInk!66, align=center] at (7.4,-2.55)
  {The label vote can be unanimous; \toolname{} trusts only evidence that is \emph{shared across agents} and \emph{re-checked on the page}.};
\end{tikzpicture}}
\caption{\toolname{} as a qualitative evidence story. Three agents inspect the
same page from complementary views, then place their cited cues on a shared
evidence board. Colored pins show provenance, while tool badges mark cues that
are confirmed or rejected by the page itself. \toolname{} acts when the shared
reasons survive these checks and returns them as an evidence packet; when agents
agree on the label but their reasons split or fail verification, it abstains.}
\Description{A qualitative left-to-right illustration of PhishProof. A suspicious
login page is inspected by three agents, their observations become evidence rows
with agent pins and tool checks, a grounded-agreement meter separates verified
shared reasons from label-only agreement, and the system either acts with an
evidence packet or abstains for review.}
\label{fig:teaser}
\end{figure*}


They share a blind spot, visible in a routine case (\autoref{fig:teaser}). A page
impersonating a bank is shown to a panel of three diverse models; all three return
``phishing,'' so every aggregation method above reports high confidence. But the
models reached that verdict for different reasons: one cited the morphed logo, one
cited a credential form posting off-domain, and one cited a certificate mismatch
that does not actually hold on the page. First, the detectors emit a label and
discard these reasons, so nothing auditable reaches the operator (\limit{1}).
Second, reliability is read at the \emph{label}, where the three agree, and is
therefore blind to the fact that they \emph{disagree on the evidence}: the verdict
is right for the wrong reason (\limit{2}). Third, even where the models do agree on
a reason, that agreement can be a shared hallucination of a cue that is not on the
page (\limit{3}). A right-label, wrong-reason verdict is not a curiosity: it is
precisely the verdict a small change to the cited cue would flip, and precisely the
one an operator should not act on blindly.

Two properties of phishing, not shared by generic question answering, make a
direct remedy possible. First, the evidence a detector relies on is
\emph{objectively checkable}: a form-action domain, a certificate organization, a
redirect target, and a rendered logo can each be re-derived from the page itself.
Second, labeled pages are cheap to assemble -- public feeds publish freshly
confirmed phishing URLs and legitimate sites supply benign pages -- so a detector's
operating point can be set on a held-out labeled split rather than guessed. The evidence that decides a verdict is also
\emph{heterogeneous}, spanning several model types (text LLMs and a
vision-language model), many signals (URL, DOM, certificate, redirect chain, logo,
and credential intent), two modalities, and a verification mechanism per signal.
Deciding which signals to gather for a given page and how to combine them is a
\emph{planning-and-aggregation} problem better suited to an agent than to a fixed
pipeline, and modern LLMs can act as such agents -- emitting typed outputs and
calling tools to gather and verify their own evidence. The opportunity is to stop
treating reliability as a property of labels and start treating it as a property of
\emph{evidence that several independent agents plan, gather, verify, and agree on,
and that the page confirms}.

To this end, we present \toolname, a multi-\emph{agent} phishing detector that
trusts a verdict in proportion to grounded cross-model agreement on the evidence.
A small, deliberately diverse panel of LLM \emph{evidence agents} investigates the
page: rather than emit a bare label, each agent \emph{plans} which signals and
modalities to examine, then gathers and grounds its evidence by calling
verification \emph{tools} -- a DOM parser, a certificate reader, a redirect tracer,
a logo matcher, and an existing reference-based or LLM brand detector -- and records
what it found in a shared, per-page set of typed \emph{evidence cues}, each carrying
the asserting agent and the confirming tool. \toolname{} \emph{aggregates} these
grounded findings: it measures agreement at the level of the cited \emph{evidence}
rather than the label, weights that agreement by whether each agreed cue verifies
against the page, and combines the two into a single Grounded Evidence-Agreement (GEA)
score. It acts only when the score is high -- returning the agreed, verified cues as
a checkable explanation -- and otherwise abstains, escalating the page to deeper
analysis or a human. A calibrator fit on a held-out labeled split maps the score to
a probability of correctness, so the operating point is set against real outcomes
rather than guessed. Reusing the existing detectors as tools also makes \toolname{}
able to match or exceed specialized detectors on full-coverage accuracy by construction,
since the panel inherits their signal. The novelty over prior multi-agent detectors is the unit of
agreement: not the label, and not one model's introspective faithfulness, but a set
of evidence cues that is simultaneously cross-agent-agreed and tool-verified.

Realizing this raises three difficulties that the design must rule out. First, to
answer \limit{2}, agreement over structured evidence must be defined so that it
cannot collapse into label agreement in disguise -- two models can name the same
brand while resting on entirely different cues. Second, to answer \limit{3}, shared hallucination must be caught
cheaply, by checking cited cues against the page rather than trusting the models to
police themselves. Cutting across both, the panel must stay small: every added
agent is another round of model and tool calls, so the gain has to come from agent
\emph{diversity} rather than count.

We make the following contributions.
\begin{compactitem}
\item \textbf{Failure mode.} We identify the \emph{right-label, wrong-reason}
failure of multi-model phishing detectors and cast trustworthy detection as
selective prediction over an evidence-agreement signal (\autoref{sec:formulation}).
\item \textbf{Method.} We introduce \toolname{}, a panel of tool-using evidence
agents that cite a shared set of typed, verified evidence cues, and its Grounded
Evidence-Agreement score, which measures cross-agent agreement over those cues and
turns it, calibrated on a held-out labeled split, into a selective act/abstain
decision with the agreed, verified cues as the explanation (\autoref{sec:method}).
\item \textbf{Evidence.} On a phishing benchmark of feed-sourced phishing pages paired
with legitimate login pages, \toolname{} matches or exceeds the re-run specialized
detectors (D1 Phishpedia, D3 PhishLLM) on full-coverage F1 while improving the
risk--coverage trade-off over single-model confidence and label voting:
it has, by a clear and statistically significant margin, the lowest expected calibration
error of any method ($0.015$) and the highest selective accuracy at $80\%$ coverage
($97.0\%$), so its reported probability of correctness can be acted on directly; its area under the risk--coverage curve (AURC) is
the lowest in point estimate ($38\%$ below the label-agreement baseline) though within the
confidence interval (CI) of the strongest baseline (\autoref{sec:experiments}).
\item \textbf{Analysis.} We isolate the contribution of evidence-level agreement,
tool-grounding, agent diversity, and calibration; cross-family \emph{diversity} and
calibration are the dominant factors, while on this benchmark -- where phishing pages
carry an unambiguous off-brand signal and few shared hallucinations -- grounding adds
little to the ranking and the right-label-wrong-reason effect is modest, appearing only
on the harder lookalike pages (\autoref{sec:experiments}, \autoref{sec:limitations}).
Under evidence-targeted attack, however, grounding earns its place: occluding the brand
logo collapses groundedness and the detector abstains, while a single-model confidence
baseline stays maximally confident on the now-wrong verdict, and the fail-safe holds even
against a white-box adversary that corrupts every cited cue at once (\autoref{sec:exp_adv}).
We further show the calibrated trust transfers to brands unseen during calibration
(\autoref{sec:exp_brands}), and that the parameter-free, small-model design is not beaten by a
learned aggregator or a $2$--$3\times$ larger panel -- the bottleneck is evidence elicitation,
not model scale (\autoref{sec:exp_negatives}).
\end{compactitem}

The rest of the paper is organized as follows. \autoref{sec:related} positions
\toolname{} against label-level reliability and multi-agent detection.
\autoref{sec:formulation} formalizes grounded evidence-agreement and the selective
objective. \autoref{sec:method} details the evidence agents, the cues they cite, the score,
and the calibrated decision rule, and
\autoref{sec:experiments} reports the evaluation.
